\documentclass{article}
\usepackage[T1,T2A]{fontenc}
\usepackage[utf8]{inputenc}
\usepackage[english,russian]{babel}
\usepackage{xcolor}
\definecolor{mygreen}{RGB}{0, 128, 0}
\title{Решение Задачи №46}
\author{Серёжи Фуканчика\\40 Ю класс}
\date{Май 2026}

\begin{document}

\maketitle
Интеграл
\section{Формулировка}
\[
\int \sqrt{1+x^2}x \,dx
\]

\section{Решение}
\subsection{Первая попытка}
Попробую применить метод интегрирования по частям:
\[
\int u \,dv=uv-\int v \,du
\]

По приоритету ЛОИЭ выбираю $u$ и $dv$:
\[
u=\sqrt{1+x^2}
\]
\[
dv=x \,dx
\]

Мне ещё понадобится $du$ (производная $u$) и $v$ (первообразная $dv$). Посчитаю $du$:

\[
du=\sqrt{1+x^2}'\,dx
\]
тут производная сложной функции, посчитаю её отдельно. Производная сложной считается по формуле $(f(g(x)))'=f'(g(x))\cdot{}g'(x)$. $f(x)=\sqrt{\cdots}$, а $g(x)=1+x^2$:

\[
\left({{(1+x^{2})}^{\frac{1}{2}}}\right)' = \frac{1}{2}{(1+x^2)}^{-\frac{1}{2}}\cdot{}2x=\frac{1}{2}\frac{2x}{\sqrt{1+x^2}}=\frac{x}{\sqrt{1+x^2}}
\]

Найду $v$ проинтегрировав $dv$:
\[
\int \,dv=\int x \,dx
\]
откуда
\[
v=\frac{x^2}{2}
\]

\[
(*)=uv-\int v \,du=\sqrt{1+x^2}\cdot\frac{x^2}{2}-\int \frac{x^2}{2} \frac{x}{\sqrt{1+x^2}} \,dx=\frac{x^2{}\sqrt{1+x^2}}{2}-\frac{1}{2}\int \frac{x^3}{\sqrt{1+x^2}} \,dx
\]

{\color{red}Упс! $u$ и $dv$ выбраны неправильно!} Получилось ещё хуже чем было.

\subsection{Вторая попытка}
Сделаю противоположный выбор $uv$:

\[
u=x
\]
\[
dv=\sqrt{1+x^2} \,dx
\]

Снова посчитаю $du$ и $v$:
\[
du=dx
\]
Но найти $v$ сложнее:
\[
\int \,dv=\int \sqrt{1+x^2} \,dx
\]
тут опять проблемы.

\subsection{Третья попытка}
Придётся делать по-другому. Попробую метод подстановки. Пусть

\[
u={1+x^2}
\]
её производная:
\[
du=2x\,dx
\]
откуда
\[
x\,dx=\frac{du}{2}
\]
т.е. исходный интеграл это:
\[
\int \sqrt{u} \frac{1}{2} \,du = \frac{1}{2} \int \sqrt{u} \,du
\]

Это табличный интеграл,
\[
\frac{1}{2} \int \sqrt{u} \,du = \frac{1}{2} \frac{u^{\frac{3}{2}}}{\frac{3}{2}} + C = \frac{1}{2}\cdot\frac{2u^{\frac{3}{2}}}{3} + C = \frac{u^{\frac{3}{2}}}{3} + C
\]
делаю обратную замену $u=1+x^2$, что даёт:
\[
\frac{\sqrt{{(1+x^2)}^3}}{3} + C
\]

\section{Ответ}
\[
\int \sqrt{1+x^2}x \,dx = \frac{\sqrt{{(1+x^2)}^3}}{3} + C
\]

\section{Проверка}
Возьму производную:
\[
\left(\frac{\sqrt{{(1+x^2)}^3}}{3}\right)'
\]
корень --- это степень $\frac{1}{2}$:
\[
\left( \frac{{(1+x^2)}^{\frac{3}{2}}}{3} \right)'
\]
вынесу константу $\frac{1}{3}$ за производную:
\[
\frac{1}{3}\cdot{}\left({(1+x^2)}^{\frac{3}{2}}\right)'
\]
применяю производную сложной функции $(f(g(x)))'=f'(g(x))\cdot{}g'(x)$:
\[
\frac{1}{3}\cdot{}\frac{3}{2}{(1+x^2)}^{\frac{1}{2}}\cdot{}2x
\]
тройки сокращаю:
\[
\frac{1}{2}{(1+x^2)}^{\frac{1}{2}}\cdot{}2x
\]
двойки тоже можно сократить:
\[
{(1+x^2)}^{\frac{1}{2}}\cdot{}x
\]
и перепишу степень $\frac{1}{2}$ как корень:
\[
\sqrt{1+x^2}\cdot{}x
\]

Отлично!


\end{document}
